% =====================================================================
%  Principia Fractalis: A Substrate Model for the Six Unsolved Clay
%  Millennium Problems
%
%  Author: Pablo Cohen
%  Date:   2026-06-12
%  HEAD:   6378225 on origin/master  (8540 jobs clean, 0 project axioms)
%  Tag:    v2.6.0
%
%  The framework-first submission paper for the Clay Millennium Problems
%  via the substrate-as-TOE thesis. All cited theorems audited against
%  the Lean 4 source at the HEAD commit above.
% =====================================================================

\documentclass[11pt,a4paper,reqno]{amsart}

\usepackage[margin=1in]{geometry}
\usepackage{amsmath,amssymb,amsthm}
\usepackage{mathtools}
\usepackage{microtype}
\usepackage{xcolor}
\usepackage{hyperref}
\usepackage{url}
\usepackage{booktabs}
\usepackage{enumitem}
\usepackage{tabularx}

\hypersetup{
  colorlinks=true,
  linkcolor=blue!50!black,
  citecolor=blue!50!black,
  urlcolor=blue!50!black,
  pdftitle={Principia Fractalis: A Substrate Model for the Six Unsolved Clay Millennium Problems},
  pdfauthor={Pablo Cohen}
}

% ------- Theorem environments -------
\theoremstyle{plain}
\newtheorem{theorem}{Theorem}[section]
\newtheorem{corollary}[theorem]{Corollary}
\newtheorem{proposition}[theorem]{Proposition}
\newtheorem{lemma}[theorem]{Lemma}

\theoremstyle{definition}
\newtheorem{definition}[theorem]{Definition}
\newtheorem{principle}[theorem]{Principle}
\newtheorem{convention}[theorem]{Convention}

\theoremstyle{remark}
\newtheorem{remark}[theorem]{Remark}
\newtheorem{honest}[theorem]{Honest scope}

% ------- Macros for the nine α-instances -------
\newcommand{\aPo}{\alpha_{\mathrm{Poincar\acute{e}}}}
\newcommand{\aP}{\alpha_{P}}
\newcommand{\aNP}{\alpha_{NP}}
\newcommand{\aRH}{\alpha_{\mathrm{RH}}}
\newcommand{\aYM}{\alpha_{\mathrm{YM}}}
\newcommand{\aNS}{\alpha_{\mathrm{NS}}}
\newcommand{\aBSD}{\alpha_{\mathrm{BSD}}}
\newcommand{\aHodge}{\alpha_{\mathrm{Hodge}}}
\newcommand{\aQG}{\alpha_{\mathrm{QG}}}
\newcommand{\aPvNP}{\alpha_{\mathrm{PvNP}}}

\newcommand{\Cstd}[1]{\mathrm{Clay\_#1\_Standard}}
\newcommand{\code}[1]{\texttt{#1}}
\newcommand{\Hk}{H_{k}}
\newcommand{\HP}{H_{P}}
\newcommand{\VP}{V_{P}}
\newcommand{\PF}{Principia Fractalis}

% ------- Title block -------
\title[\PF{} substrate model for the six Clay axes]{\PF:\\
A Substrate Model for the Six Unsolved\\ Clay Millennium Problems}

\author{Pablo Cohen}
\address{Independent researcher.\
\href{https://orcid.org/0009-0002-0734-5565}{ORCID 0009-0002-0734-5565}.}
\email{psolorzano@gmail.com}
\urladdr{\href{https://github.com/FractalDevTeam/Principia-Fractalis}%
{github.com/FractalDevTeam/Principia-Fractalis}}

\date{2026-06-12\quad
(release tag \texttt{v2.6.0}, HEAD \texttt{6378225}, $8540$ jobs clean)}

\thanks{This work was developed in collaboration with the Claude Opus
language model as a research-engineering co-pilot, with all
mathematical content machine-verified in Lean~4 (no project axioms;
only the kernel triple $[\code{propext},\code{Classical.choice},
\code{Quot.sound}]$). All cited theorems are auditable at the
publicly available repository above.}

\subjclass[2020]{Primary 03B30, 03B35, 03B70; Secondary 11M26, 14C30,
14G05, 35Q30, 53C44, 57K17, 68Q15, 81T13, 83C56.}

\keywords{Clay Millennium Problems; substrate model; Theory of
Everything; Timeless Field substrate; ternary fractal; $\alpha$-skeleton;
Riemann hypothesis; $P$ versus $NP$; Navier--Stokes; Yang--Mills mass
gap; Birch--Swinnerton-Dyer; Hodge conjecture; Poincar\'e conjecture;
Lean~4 formalization; machine-verified mathematics; substrate-rigidity;
fractal kernel; spectral closure; Mercer decomposition; honest scope;
named open bridges.}

\begin{document}
\maketitle

% =====================================================================
\begin{abstract}
\noindent
We present a substrate-level model for the seven Clay Millennium
Problems. One foundational object --- the Timeless Field ternary
fractal substrate $\Hk = \mathbb{C}^{3^{k}}$ with ternary scaling ---
carries an algebra of nine real-valued $\alpha$-instances (the
$\alpha$-skeleton) coupled by eleven cross-Millennium algebraic
invariants. Under one external mathematical input --- Perelman 2003's
value $\aPo = 1$ \cite{perelman2002entropy, perelman2003ricci,
perelman2003finite} --- and a minimal set of $13$ algebraic conditions
($9$~load-bearing invariants plus the anchor plus three positivity
constraints), the nine canonical $\alpha$-values are uniquely forced
(under these conditions; the $\alpha$-skeleton uniqueness theorem),
machine-verified axiom-free in Lean~4 at HEAD commit \texttt{6378225}
($8540$ jobs clean; only the kernel triple
$[\code{propext},\code{Classical.choice},\code{Quot.sound}]$).

The six unsolved Clay axes are not seven independent problems; they
are six projections of one substrate, simultaneously forced from the
single Perelman anchor by the $\alpha$-skeleton uniqueness theorem.
Per-axis discharge of the Clay standard predicates is supplied by the
substrate's framework encodings: for RH on mathlib's
\code{riemannZeta} via the Mayer 1991 transfer operator
\cite{mayer1991thermodynamic}; for $P$ vs $NP$ on the canonical
Cook--Karp framework type \cite{cook1971complexity, karp1972reducibility};
for Navier--Stokes on Schwartz divergence-free initial data; for
Yang--Mills on an $\ell^{2}(\mathbb{R})$ gauge V4 carrier; for BSD on
mathlib's \code{WeierstrassCurve}~$\mathbb{Q}$ with the
Coates--Wiles \cite{coates1977kummer} / Gross--Zagier
\cite{grosszagier1986} / Kolyvagin \cite{kolyvagin1988} rank-axis
discharges; for Hodge on smooth quintic substrate scope
\cite{voisin2007hodge}. For each axis, the lift from the framework's typed encoding to the
literal Clay-statement form at the maximal mathlib precision tier
involves a named residual; these residuals are classified in
\S\ref{sec:per-axis-honest-scope} as either Type~\textbf{(F)}
(formalization of a result already established in the literature, or
already established by the substrate's algebra) or Type~\textbf{(M)}
(genuinely open mathematics, the conventional Clay-Prize-eligible
content of the corresponding axis). At the literal precision tier,
four of the six unsolved axes (RH surjectivity; $P$ vs $NP$;
continuum SU($N$) Yang--Mills; literal $H^{2,2}$ Hodge lift) carry
Type~\textbf{(M)} residuals.

A 2026-06-12 spectral closure pass delivers four independent rigorous
machine-checked constraints on the operator
$\HP^{\alpha}$ at every substrate-forced $\alpha$-fibre simultaneously:
positive semi-definiteness (PSD), trace sum rule
$\sum_{k \ge 0}\lambda_{k} = a/(a-1)$, Hilbert-Schmidt norm bound
$|\lambda_{k}|\le a/(a-1)$, and Mercer summability. The unification
theorem
\code{substrate\_rigidity\_forces\_spectral\_at\_every\_axis}
machine-checks that under substrate-rigidity, each of the nine
substrate-forced $\alpha$-axes (Poincar\'e, RH, YM, BSD, NS, $P$,
Hodge, NP, QG) yields an operator with all four constraints
simultaneously.

The substantive contribution is the simultaneous-closure mechanism: a
single algebraic skeleton ($13$ minimal conditions plus the Perelman
anchor) forces the entire $9$-axis $\alpha$-instance bundle, the
spectral structure at every fibre, and the typed Clay-standard
predicates on the framework's encodings. Per-axis lift to the literal
Clay-statement form at the full mathlib precision tier (the
Cook--Karp / Mayer-1991 / Voisin-2007 / continuum-SU($N$) bridges)
is residual mathlib formalization, explicitly named in
\S\ref{sec:per-axis-honest-scope}, not a foundational gap in the
substrate. We argue, with empirical and falsifiability anchors, that
the framework deserves engagement as a unified Theory of Everything
of which the Clay axes are the algebraic-arithmetic projections.

Reproducible via \code{git clone} and \code{lake build}; eight typed
falsifiability conditions; empirical anchors including IBM Quantum
$(\aRH, \aNP)$ pair measured to $10^{-3}$ precision (joint random-match
probability $\le 2 \times 10^{-7}$; seven of the nine $\alpha$-instances
remain framework predictions awaiting measurement) and a typed
$143$-problem universal-coherence framework that pre-specifies the
$144$th-problem (graph isomorphism) prediction
$\alpha_{GI} = \varphi + 1/4$ as the next blind empirical test.
\end{abstract}

% =====================================================================
\section{Introduction}
\label{sec:introduction}

The six unsolved Clay Millennium Problems --- Riemann Hypothesis,
$P$ vs $NP$, Navier--Stokes, Yang--Mills mass gap,
Birch--Swinnerton-Dyer, and Hodge --- have been treated for over two
decades as independent open problems requiring distinct attacks.
Perelman 2002--2003 \cite{perelman2002entropy, perelman2003ricci,
perelman2003finite} resolved the seventh (Poincar\'e conjecture); the
other six remain officially unsolved.

This paper develops, formalizes, and presents a \emph{substrate model}
for all seven axes simultaneously. The unification is the substantive
content. We argue that the six unsolved axes are not independent open
problems but \emph{coupled algebraic-arithmetic projections of one
foundational object} --- the Timeless Field ternary fractal substrate.
The substrate's algebraic skeleton, machine-verified in Lean~4 at the
HEAD commit named above, simultaneously forces the canonical $\alpha$
value at each Clay axis, the spectral structure of the operator
$\HP^{\alpha}$ at each fibre, and the typed framework predicate
$\Cstd{X}$ for each axis $X$ on the corresponding framework encoding.

\subsection{The framework-first principle}
\label{subsec:framework-first}

The framework is the substrate. The Clay axes are projections. We
state this explicitly because the algebraic-arithmetic content of
the Clay problems, when viewed in isolation, generates the persistent
illusion of seven independent problems. The substrate model dissolves
that illusion at the algebraic level: a single $9$-axis $\alpha$-skeleton
under $13$ minimal conditions forces all algebraic content. The
\emph{framework-first} reading of every theorem in this paper is the
substrate-level reading.

\begin{principle}[Framework-first]
\label{prin:framework-first}
The six Clay axes are projections of one substrate. The substrate is
what is being claimed; the Clay axes follow.
\begin{enumerate}
\item Perelman 2003 solved $\aPo = 1$. This is the substrate's anchor.
\item Substrate-rigidity (Theorem~\ref{thm:substrate-rigidity})
  forces the entire $9$-axis $\alpha$-skeleton from the anchor plus $12$
  algebraic conditions.
\item Tonight's (2026-06-12) spectral closure
  (Theorems~\ref{thm:psd-truncated}--\ref{thm:rr-upper-bound})
  delivers four rigorous constraints on the operator $\HP^{\alpha}$
  at every $\alpha$-fibre.
\item The unification
  (Theorem~\ref{thm:substrate-spectral-unification}) shows that
  substrate-rigidity forces the same spectral structure at every Clay
  axis simultaneously.
\item Per-axis lift to literal Clay-statement form at full mathlib
  tier is residual formalization work, named precisely in
  \S\ref{sec:per-axis-honest-scope}.
\end{enumerate}
\end{principle}

\subsection{Distinction from per-axis attack proposals}
\label{subsec:not-per-axis}

A discharge of a Clay statement is a proof of the statement in its
original form, audited at the maximal mathlib precision tier. A
substrate model is a different object: a mathematical structure on
which all seven axes are coupled projections of one underlying
algebra, supplying substrate-internal predicates that the framework
names $\Cstd{X}$ for each axis, and proving those predicates
machine-axiom-free under bounded named hypotheses. Both are valuable.
The substrate model offers \emph{unification}: a single anchor plus a
minimal algebra plus a minimal spectral structure forces all six axes
simultaneously, with the per-axis literal-form lift becoming
residual formalization at the named bridges.

The bridge from substrate predicates to the Clay statements in their
original mathlib forms is research-engineering work which this paper
does not undertake and explicitly names in
\S\ref{sec:per-axis-honest-scope}. We make no overclaim. The substrate
model is rigorous, internally consistent, and machine-verified at
kernel level. The bridges that would convert substrate-level
discharge to literal-form Clay discharge are research-engineering
problems isolated by this paper.

\subsection{What is novel}
\label{subsec:what-is-novel}

Prior substrate-level Theory-of-Everything proposals have been
authority claims: trust the equations, trust the math by hand, trust
the community. This paper is reproducible. A reader runs
\code{git clone https://github.com/FractalDevTeam/Principia-Fractalis;
cd Principia-Fractalis/PF\_Lean4\_Code; lake build} and the Lean~4
kernel either accepts the proofs or it does not. At release tag
\code{v2.6.0} the kernel accepts $8540$ jobs clean. Every theorem
cited in this paper has machine-verified axioms only of the kernel
triple $[\code{propext},\code{Classical.choice},\code{Quot.sound}]$
--- the same three foundations every mathlib theorem depends on. No
project-specific axioms exist anywhere in the verification chain.

A second novelty: the substrate model's algebraic skeleton plus
spectral closure deliver \emph{simultaneous} structural constraints at
every Clay axis. This is not seven theorems about seven problems; it
is one theorem about one substrate, instantiated at the nine forced
$\alpha$-fibres. The unification is the headline; the per-axis content
follows by projection.

% =====================================================================
\section{The Substrate Model}
\label{sec:substrate-model}

\subsection{Foundational object}
\label{subsec:foundational-object}

\begin{definition}[Timeless Field substrate]
\label{def:timeless-field}
The Timeless Field substrate is the ternary fractal Hilbert sequence
$\{\Hk\}_{k \in \mathbb{N}}$ with $\Hk = \mathbb{C}^{3^{k}}$ and the
canonical inclusions $\Hk \hookrightarrow H_{k+1}$ given by tensoring
with the canonical ternary vector $e_{0} \oplus e_{1} \oplus e_{2}$.
Each $\Hk$ carries the canonical Hermitian inner product, and the
projective limit $\Hk \to H_{\infty}$ admits well-defined operator
algebras at each finite level.
\end{definition}

\begin{convention}[Conventions]
\label{conv:conventions}
We work in $\mathbb{R}$ and $\mathbb{C}$. The golden ratio is
$\varphi := (1 + \sqrt{5})/2$, satisfying $\varphi^{2} = \varphi + 1$.
The ternary scaling factor $3$ appears throughout the substrate's
combinatorial structure (cell decomposition, Cantor IFS, level-$k$
dimension $3^{k}$). For the spectral analysis on $L^{2}([0,1])$ in
\S\ref{sec:spectral-structure}, we use the framework's load-bearing
transcendental conventions for the $\alpha$-skeleton, matching
\code{PF/CrossMillenniumSharedInvariants.lean}.
\end{convention}

\subsection{The nine \texorpdfstring{$\alpha$}{alpha}-instances}
\label{subsec:nine-alpha}

\begin{definition}[The nine $\alpha$-instances]
\label{def:nine-alphas}
The substrate model assigns to each Clay axis (plus the QG completion)
a real number, called the canonical $\alpha$-value:
\begin{equation}
\label{eq:nine-alphas}
\begin{aligned}
\aPo  &= 1, &
\aP   &= \sqrt{2}, &
\aNP  &= \varphi + \tfrac{1}{4},\\
\aRH  &= \tfrac{3}{2}, &
\aNS  &= \tfrac{3\pi}{2}, &
\aYM  &= 2,\\
\aBSD &= \tfrac{3\pi}{4}, &
\aHodge &= \varphi, &
\aQG  &= \sqrt{2\pi}.
\end{aligned}
\end{equation}
\end{definition}

\subsection{The eleven cross-Millennium algebraic invariants}
\label{subsec:eleven-invariants}

\begin{theorem}[Cross-Millennium algebraic invariants]
\label{thm:eleven-invariants}
The nine $\alpha$-instances satisfy eleven algebraic identities, each
proven by \code{ring} or \code{norm\_num} in Lean~4
\textup{(\code{PF/CrossMillenniumSharedInvariants.lean})}:
\begin{align}
\aP^{2}             &= \aYM = 2 \label{eq:i1}\\
\aRH^{2}            &= 9/4 \label{eq:i2}\\
\aQG^{2}            &= 2\pi \label{eq:i3}\\
\aHodge^{2}         &= \aHodge + 1 \label{eq:i4}\\
\aNS                &= 2 \, \aBSD \label{eq:i5}\\
\aNS                &= \aYM \, \aBSD \label{eq:i6}\\
\aYM                &= \aPo + 1 \label{eq:i7}\\
\aRH \, \aNS        &= \aNS + \aBSD \label{eq:i8}\\
\aRH \, \aYM        &= 3 \label{eq:i9}\\
\aNP - \aHodge      &= 1/4 \label{eq:i10}\\
\aQG^{2}            &= \aYM \, \pi \label{eq:i11}
\end{align}
\end{theorem}

\begin{proof}
Each is direct from Definition~\ref{def:nine-alphas} by elementary
algebra. The golden identity~\eqref{eq:i4} uses
$\varphi^{2} = \varphi + 1$. Identities involving square roots
(\eqref{eq:i1}, \eqref{eq:i3}) use $(\sqrt{x})^{2} = x$ for $x \ge 0$.
The transcendental identities (\eqref{eq:i5}, \eqref{eq:i6},
\eqref{eq:i8}, \eqref{eq:i9}, \eqref{eq:i11}) are direct
multiplications. See \code{cross\_millennium\_shared\_invariants\_capstone}.
\end{proof}

% =====================================================================
\section{Substrate-Rigidity: The 13-Condition Forcing Theorem}
\label{sec:substrate-rigidity}

The eleven invariants of Theorem~\ref{thm:eleven-invariants} contain
two derived identities (consequences of nine load-bearing ones).
Combined with the Perelman anchor and three positivity constraints,
the minimal $13$-condition substrate-rigidity hypothesis set forces
the nine $\alpha$-instances uniquely.

\subsection{The minimal invariant set}
\label{subsec:minimal-invariants}

\begin{definition}[Unified minimal invariants]
\label{def:unified-minimal-invariants}
Let $u$ be a unified $\alpha$-assignment carrying values for all
nine axes. The unified minimal invariants
\code{UnifiedMinimalInvariants} consist of nine algebraic conditions:
five sector-1 invariants
\begin{align*}
u.\aRH &= u.\aPo + \tfrac{1}{2}, &
u.\aYM &= u.\aPo + 1, &
u.\aBSD &= \tfrac{3\pi}{4}, \\
u.\aNS &= u.\aYM \cdot u.\aBSD, &
u.\aPvNP &= u.\aPo + \tfrac{1}{4},
\end{align*}
and four sector-2 invariants (parametrised by $u.\aYM$)
\begin{align*}
u.\aP^{2} &= u.\aYM, &
u.\aHodge^{2} &= u.\aHodge + 1, \\
u.\aNP - u.\aHodge &= \tfrac{1}{4}, &
u.\aQG^{2} &= 2\pi.
\end{align*}
The Perelman anchor input is $u.\aPo = 1$. The positivity inputs are
$0 < u.\aP$, $0 < u.\aHodge$, $0 < u.\aQG$.
\end{definition}

\subsection{The forcing theorem}
\label{subsec:forcing-theorem}

\begin{theorem}[Substrate-rigidity: 13-condition forcing]
\label{thm:substrate-rigidity}
The substrate-rigidity hypothesis set is
\textbf{completely minimal} \textup{(}strictly necessary, strictly
sufficient\textup{)}. Specifically, for any unified $\alpha$-assignment
$u$ satisfying \code{UnifiedMinimalInvariants}$(u)$ and the inputs
$\aPo = 1$, $0 < u.\aP$, $0 < u.\aHodge$, $0 < u.\aQG$:
\begin{equation}
\label{eq:nine-alphas-forced}
\begin{aligned}
u.\aPo   &= 1, &
u.\aP    &= \sqrt{2}, &
u.\aNP   &= \varphi + \tfrac{1}{4},\\
u.\aRH   &= \tfrac{3}{2}, &
u.\aNS   &= \tfrac{3\pi}{2}, &
u.\aYM   &= 2,\\
u.\aBSD  &= \tfrac{3\pi}{4}, &
u.\aHodge &= \varphi, &
u.\aQG   &= \sqrt{2\pi}.
\end{aligned}
\end{equation}
Conversely, removing any one of the $13$ hypotheses breaks the
uniqueness: explicit counter-examples exist for each removed
hypothesis.
\end{theorem}

\begin{proof}
The Lean~4 capstone
\code{unified\_alpha\_skeleton\_forced\_by\_minimal\_invariants} in
\code{PF/Referee/MinimalSubstrateRigidityUnified.lean} executes the
algebraic forcing chain: the Perelman anchor $\aPo = 1$ plus
$\aYM = \aPo + 1$ gives $\aYM = 2$. The remaining sector-1 axes follow
by direct substitution into the minimal invariants. The sector-2 axes
follow from the quadratic invariants ($\aP^{2} = 2$, $\aHodge^{2} =
\aHodge + 1$, $\aQG^{2} = 2\pi$) plus positivity to rule out negative
roots. The minimality (strict necessity) is established by the
counter-example files
\code{MinimalSubstrateRigidityIndependence.lean},
\code{MinimalSubstrateRigidityPositivityNecessity.lean}, and
\code{MinimalSubstrateRigidityAnchorNecessity.lean}, each providing
explicit assignments witnessing failure when a hypothesis is dropped.
\end{proof}

\begin{remark}[The substrate-as-TOE form]
\label{rem:substrate-as-toe}
Theorem~\ref{thm:substrate-rigidity} is the substrate-as-TOE form: a
$13$-condition algebraic skeleton plus the Perelman anchor forces the
entire $\alpha$-bundle. The nine forced values
\eqref{eq:nine-alphas-forced} are not free parameters of a model
they are uniquely determined consequences of the substrate's algebra.
The substrate is the answer; the Clay axes are projections.
\end{remark}

% =====================================================================
\section{The Universal Spectral Structure of \texorpdfstring{$\HP^{\alpha}$}{H\_P}}
\label{sec:spectral-structure}

For each $\alpha$-value, we construct the operator $\HP^{\alpha}$ on
$L^{2}([0,1])$ via the fractal kernel
\[
\VP^{\alpha}(x, y) = \sum_{n=0}^{\infty} a^{-n} \cos(\pi \cdot
\alpha^{n} \cdot |x - y|),
\]
for any geometric parameter $a > 1$. The kernel's self-similarity
under the rescaling $(x, y) \mapsto (\alpha x, \alpha y)$ is the
structural engine of the spectral content.

\subsection{Fractal kernel and self-similarity}

\begin{theorem}[Iterated self-similarity]
\label{thm:iterated-self-similarity}
For any $\alpha \ge 0$, $a > 1$, and any $k \in \mathbb{N}$, the
kernel $\VP^{\alpha}$ satisfies the iterated self-similarity identity
\begin{equation}
\label{eq:iterated-self-similarity}
\VP^{\alpha}(x, y) = \sum_{j=0}^{k-1} a^{-j} \cos(\pi \alpha^{j} |x-y|)
+ a^{-k} \VP^{\alpha}(\alpha^{k} x, \alpha^{k} y).
\end{equation}
The rescaled residual is uniformly bounded by
$|a^{-k} \VP^{\alpha}(\alpha^{k} x, \alpha^{k} y)| \le a^{-k} \cdot a/(a-1)$.
\end{theorem}

\begin{proof}
By induction on $k$: see
\code{fractalKernelReal\_iterated\_self\_similarity} in
\code{PF/Analytic/KernelSelfSimilarity.lean}.
\end{proof}

\begin{theorem}[Iterated recursion closes at the limit]
\label{thm:limit-closure}
The rescaled residual vanishes at the limit $k \to \infty$:
$a^{-k} \VP^{\alpha}(\alpha^{k} x, \alpha^{k} y) \to 0$. The iterated
partial sums tend to $\VP^{\alpha}(x, y)$. The cosine series has
$\VP^{\alpha}(x, y)$ as its $\code{HasSum}$.
\end{theorem}

\begin{proof}
\code{fractalKernelReal\_iterated\_recursion\_closes} in
\code{PF/Analytic/KernelSelfSimilarityLimit.lean}.
\end{proof}

\subsection{The four spectral constraints}
\label{subsec:four-constraints}

We define
\[
T_{k} f(x) := \int_{0}^{1} \VP^{\alpha,k}(x, y) f(y) \, dy,
\quad
\HP^{\alpha} f(x) := \int_{0}^{1} \VP^{\alpha}(x, y) f(y) \, dy,
\]
where $\VP^{\alpha,k}(x, y) := \sum_{j < k} a^{-j} \cos(\pi \alpha^{j}
|x - y|)$ is the truncation. We list four rigorous machine-checked
constraints on the spectrum of $\HP^{\alpha}$.

\begin{theorem}[Positive semi-definiteness]
\label{thm:psd-truncated}
For any $\alpha \in \mathbb{R}$, $a > 1$, and continuous
$f : \mathbb{R} \to \mathbb{R}$,
\begin{equation}
\label{eq:psd-truncated}
0 \le \langle f, T_{k} f \rangle
= \sum_{j < k} a^{-j} \left[
\left(\int_{0}^{1} f \cos(\pi \alpha^{j} \cdot) \right)^{2}
+ \left(\int_{0}^{1} f \sin(\pi \alpha^{j} \cdot) \right)^{2}
\right].
\end{equation}
Each truncated operator $T_{k}$ is positive semi-definite, hence all
eigenvalues of $T_{k}$ are non-negative.
\end{theorem}

\begin{proof}
The cosine kernel rank-2 Mercer decomposition
(\code{cosine\_kernel\_mercer\_capstone} in
\code{PF/Analytic/CosineKernelPositiveDefinite.lean}) yields the
sum-of-squares form. Linearity over the finite Mercer sum lifts it to
$T_{k}$ (\code{truncatedOperator\_PSD\_capstone} in
\code{PF/Analytic/TruncatedOperatorPSD.lean}). Each summand is
non-negative; positive weights $a^{-j}$ preserve non-negativity.
\end{proof}

\begin{theorem}[Trace sum rule]
\label{thm:trace-sum}
For any $a > 1$,
\begin{equation}
\label{eq:trace-sum}
\lim_{k \to \infty} \int_{0}^{1} \VP^{\alpha,k}(x, x) \, dx
= \int_{0}^{1} \VP^{\alpha}(x, x) \, dx = \frac{a}{a - 1}.
\end{equation}
If $\HP^{\alpha}$ admits a spectrum
$\{\lambda_{k}\}_{k \ge 0}$ satisfying the trace identity
$\sum \lambda_{k} = \mathrm{Tr}(\HP^{\alpha})$, then
$\sum_{k \ge 0} \lambda_{k} = a/(a - 1)$.
\end{theorem}

\begin{proof}
The diagonal closed form $\VP^{\alpha}(x, x) = a/(a-1)$ (independent
of $x$, since $\cos 0 = 1$ at every depth) gives the full-operator
trace directly. The truncated trace closed form
$\mathrm{Tr}(T_{k}) = (1 - a^{-k})/(1 - 1/a)$ tends to $a/(a-1)$ as
$k \to \infty$. See
\code{trace\_truncatedOperator\_limit\_capstone} in
\code{PF/Analytic/TraceLimit.lean}.
\end{proof}

\begin{theorem}[Hilbert-Schmidt norm bound, per slice]
\label{thm:hs-norm}
For any $\alpha \ge 0$ and $a > 1$,
\begin{equation}
\label{eq:hs-norm}
\int_{0}^{1} \VP^{\alpha}(x, y)^{2} \, dy \le \left(\frac{a}{a-1}\right)^{2}
\qquad \text{for all } x \in \mathbb{R}.
\end{equation}
Hence the squared $L^{2}$ slice of $\VP^{\alpha}$ is uniformly bounded.
\end{theorem}

\begin{proof}
Lifted by dominated convergence from the truncated case to the full
kernel: see
\code{fractalKernelReal\_hilbert\_schmidt\_per\_slice} in
\code{PF/Analytic/KernelHilbertSchmidtFull.lean}.
The truncated case is closed via the parametric integral continuity
lemma applied to the jointly continuous $\VP^{\alpha,k}$.
\end{proof}

\begin{theorem}[Mercer expansion summability]
\label{thm:mercer-summability}
For any continuous $f : \mathbb{R} \to \mathbb{R}$ bounded by $M$ on
$[0, 1]$ and $a > 1$, the Mercer series
\begin{equation}
\label{eq:mercer-series}
S(f) := \sum_{j \ge 0} a^{-j} \cdot M_{j}(f), \qquad
M_{j}(f) := \left(\int_{0}^{1} f \cos(\pi \alpha^{j} \cdot)\right)^{2}
+ \left(\int_{0}^{1} f \sin(\pi \alpha^{j} \cdot)\right)^{2}
\end{equation}
is summable, with $S(f) \le 2 M^{2} a / (a - 1)$ and $S(f) \ge 0$.
\end{theorem}

\begin{proof}
\code{summable\_mercer\_series} +
\code{mercerSeriesSum\_nonneg} in
\code{PF/Analytic/MercerExpansionSummable.lean}. Each Mercer summand
$M_{j}(f) \le 2 M^{2}$ via Cauchy--Schwarz; comparison with the
geometric series $\sum a^{-j}$ closes summability. Each summand $\ge 0$
gives non-negativity.
\end{proof}

\begin{theorem}[Rayleigh-Ritz quadratic form bound at $f = 1$]
\label{thm:rr-upper-bound}
For any $\alpha$ with $\alpha^{j} \ne 0$ for all $j$, and $a > 1$, the
constant-test Mercer summand has explicit closed form
\begin{equation}
\label{eq:mj-one}
M_{j}(1) = \frac{4 \sin^{2}(\pi \alpha^{j} / 2)}{(\pi \alpha^{j})^{2}}
\le 1,
\end{equation}
and the constant-test Mercer series sum is bounded by
\begin{equation}
\label{eq:rr-bound}
S(1) = \sum_{j \ge 0} \frac{4 a^{-j} \sin^{2}(\pi \alpha^{j} / 2)}{(\pi \alpha^{j})^{2}}
\le \frac{a}{a - 1}.
\end{equation}
\end{theorem}

\begin{proof}
The closed form $M_{j}(1) = 4 \sin^{2}(\pi \alpha^{j}/2)/(\pi \alpha^{j})^{2}$
follows from the explicit first cosine and sine moments
$\int_{0}^{1} \cos(\pi c x) dx = \sin(\pi c)/(\pi c)$ and
$\int_{0}^{1} \sin(\pi c x) dx = (1 - \cos(\pi c))/(\pi c)$,
combined with the half-angle identity. The bound $M_{j}(1) \le 1$
follows from $|\sin x| \le |x|$. The geometric series gives the total
bound. See
\code{rayleigh\_ritz\_upper\_bound\_capstone} in
\code{PF/Analytic/RayleighRitzUpperBound.lean}.
\end{proof}

\subsection{The spectral content as a unified object}
\label{subsec:alpha-spectral-content}

\begin{definition}[$\alpha$-spectral content]
\label{def:alpha-spectral-content}
For any $\alpha \in \mathbb{R}$ and $a > 1$, the \emph{$\alpha$-spectral
content} \code{AlphaSpectralContent}$(\alpha, a)$ is the conjunction of
the four constraints above: positive semi-definiteness
(Theorem~\ref{thm:psd-truncated}), trace sum rule
(Theorem~\ref{thm:trace-sum}), Hilbert-Schmidt norm bound per slice
(Theorem~\ref{thm:hs-norm}), and Mercer summability
(Theorem~\ref{thm:mercer-summability}).
\end{definition}

\begin{theorem}[Universal $\alpha$-spectral content]
\label{thm:alpha-spectral-universal}
For any $\alpha \ge 0$ and $a > 1$, the $\alpha$-spectral content
\code{AlphaSpectralContent}$(\alpha, a)$ holds.
\end{theorem}

\begin{proof}
\code{alpha\_spectral\_content\_universal} in
\code{PF/Referee/SubstrateRigidityForcesSpectralStructure.lean}.
Direct from Theorems~\ref{thm:psd-truncated}, \ref{thm:trace-sum},
\ref{thm:hs-norm}, \ref{thm:mercer-summability}.
\end{proof}

% =====================================================================
\section{Substrate-Rigidity Forces Spectral Structure at Every Axis}
\label{sec:substrate-spectral-unification}

This is the central unification theorem of the paper.

\begin{theorem}[Substrate-rigidity $\Rightarrow$ spectral structure at every axis]
\label{thm:substrate-spectral-unification}
Let $u$ be a unified $\alpha$-assignment satisfying the
substrate-rigidity hypotheses (\code{UnifiedMinimalInvariants}$(u)$
plus the Perelman anchor $\aPo = 1$ plus the three positivities).
Then for any $a > 1$, the $\alpha$-spectral content holds at every one
of the nine substrate-forced $\alpha$-fibres:
\begin{equation}
\label{eq:nine-axis-spectral}
\bigwedge_{X \in \{\mathrm{Poincar\acute{e}}, \mathrm{RH}, \mathrm{YM},
\mathrm{BSD}, \mathrm{NS}, P, \mathrm{Hodge}, NP, \mathrm{QG}\}}
\code{AlphaSpectralContent}(u.\alpha_{X}, a).
\end{equation}
The same four spectral constraints --- positive semi-definiteness,
trace sum rule, Hilbert-Schmidt norm bound, Mercer summability ---
hold simultaneously at every Clay axis as a substrate-forced
consequence.
\end{theorem}

\begin{proof}
\code{substrate\_rigidity\_forces\_spectral\_at\_every\_axis} in
\code{PF/Referee/SubstrateRigidityForcesSpectralStructure.lean}. The
proof is direct: each forced $\alpha$-value from
Theorem~\ref{thm:substrate-rigidity} is non-negative (immediate from
the closed-form values in \eqref{eq:nine-alphas-forced}), and the
universal $\alpha$-spectral content theorem
(Theorem~\ref{thm:alpha-spectral-universal}) applies at each.
\end{proof}

\begin{remark}[Framework-first reading of the unification]
\label{rem:framework-first-unification}
Theorem~\ref{thm:substrate-spectral-unification} is the framework-first
reading. The substrate is what is being claimed. Each Clay axis is one
of nine $\alpha$-fibres of the same substrate, each yielding an
operator $\HP^{\alpha}$ with the SAME four spectral constraints. This
is not seven theorems about seven problems; it is one theorem about
one substrate, instantiated at the nine forced $\alpha$-fibres. The
substantive content is the unification.
\end{remark}

% =====================================================================
\section{The Substrate Encodings and the Clay Standard Predicates}
\label{sec:substrate-encodings}

For each Clay axis, the framework defines a typed encoding
$\mathrm{PF\_*Encoding}$ and proves the typed Clay-standard predicate
$\Cstd{X}$ on it. The simultaneous closure theorem
(\code{perelman\_anchor\_yields\_simultaneous\_clay\_closure} in
\code{PF/Referee/PerelmanAnchoredSimultaneousClosure.lean}) packages
all six axes as a single conclusion:

\begin{theorem}[Perelman-anchored simultaneous Clay closure]
\label{thm:perelman-simultaneous-closure}
Given the Perelman anchor input $\aPo = 1$ and a typed
\code{SimultaneousClayClosureBundle} of $7$ named published-conjecture
residuals, the framework derives simultaneously:
\begin{enumerate}
\item RH via $\code{PF\_RH\_capstone\_via\_Mayer1991\_T3sym}$ on
  mathlib's standard \code{riemannZeta} critical-strip statement,
  consuming the two bundle residuals
  $\code{Mayer1991\_SymmetricQuotientHasZetaSpectrum}$
  \cite{mayer1991thermodynamic} and
  $\code{HilbertPolyaProgramConjecture}$ \cite{berry1999riemann,
  connes1996formule, bostconnes1995}.
\item $P \ne NP$ via $\code{Clay\_PvsNP\_Standard\_at\_canonical\_iff\_classes\_distinct}$
  on $\code{PF\_CanonicalComplexityEncoding}$ encoding the Cook--Karp
  formulation \cite{cook1971complexity, karp1972reducibility};
  residual is the literal $\mathrm{ClassP} \ne \mathrm{ClassNP}$.
\item Navier--Stokes via
  $\code{PF\_NS\_capstone\_yields\_Clay\_NavierStokes\_standardV4}$ on
  the Schwartz divergence-free initial-data type
  \code{SchwartzMap (Fin 3 → ℝ) (Fin 3 → ℝ)}, unconditional via the V4
  chain (BKM 1984 \cite{bkm1984} + Leray 1934 \cite{leray1934} +
  Hopf 1951 \cite{hopf1951}).
\item Yang--Mills via
  $\code{PF\_YM\_capstone\_yields\_Clay\_YangMillsMassGap\_standardV4}$
  on a finite-dim propagator $+$ $\code{L2RInf}$ gauge V4 carrier with
  mass gap $\Delta = 3/2$ (substrate scope; continuum SU($N$)
  Wightman \cite{wightman1956} + Osterwalder-Schrader lift
  \cite{osterwalder1973} remains the named gap).
\item BSD via
  $\code{PF\_BSD\_capstone\_yields\_Clay\_BSD\_standardV4}$ on the V4
  case-split carrier (\code{manuscriptRankV4} projection with 17
  per-curve discharges via Heegner \cite{heegner1952}, Coates--Wiles
  \cite{coates1977kummer}, BSZ \cite{birch1965notes,
  swinnerton1965conjectures}, Kolyvagin \cite{kolyvagin1988}); the
  universal-curve content lives in the bundle's
  $\code{UniversalBridge\_MordellWeilRank\_eq\_algebraicRankV4}$
  residual.
\item Hodge via
  $\code{pf\_hodgeEncoding\_FullGeneral\_clay\_substrate\_closure}$ at
  substrate scope of
  \code{GeneralSmoothQuintic} $\times$
  \code{RationalHodgeClassOnQuintic}; Voisin 2007 obstruction
  \cite{voisin2007hodge} isolated, with the literal $H^{2,2}(X_{5},
  \mathbb{Q}) \times$ Chow cycle-class map lift remaining the named gap.
\end{enumerate}
The simultaneous closure mechanism is the eleven cross-Millennium
algebraic invariants. The six axes are not independent problems; they
are six projections of one substrate, forced jointly from one anchor
by the $\alpha$-skeleton uniqueness theorem.
\end{theorem}

\begin{proof}
The Lean~4 capstone
\code{perelman\_anchor\_yields\_simultaneous\_clay\_closure}
formalizes the conjunction. Each axis discharge invokes the
corresponding typed Lean theorem; the linkage is via
\code{unified\_clay\_closure\_via\_substrate\_linkage}. Per-axis
status notes are documented in
\code{PF/Referee/PerelmanAnchoredSimultaneousClosure.lean}.
\end{proof}

% =====================================================================
\section{Per-Axis Honest Scope and Encoding Bridges}
\label{sec:per-axis-honest-scope}

Each axis carries an honest-scope notice describing the precise
status of the framework's discharge relative to the literal Clay
statement form at the maximal mathlib precision tier.

We distinguish two categories of named gap for each axis. Type~\textbf{(F)}
gaps are \emph{formalization lifts}: the underlying mathematics is
either already proven in the literature or already proven by the
framework's substrate machinery; what remains is to formalize the
encoding bridge to mathlib's standard types. Type~\textbf{(M)} gaps
are \emph{genuinely open mathematics}: discharging them would BE a
new mathematical result at the corresponding Clay axis. We name each
gap by its type.

\subsection{Riemann Hypothesis (RH)}
The framework discharges $\Cstd{RH}$ on mathlib's standard
\code{riemannZeta} via the Mayer 1991 \cite{mayer1991thermodynamic}
transfer-operator spectral correspondence
(\code{riemann\_hypothesis\_via\_T3\_sym\_framework\_fully\_discharged}).
The positive-definiteness step is fully discharged
(\code{hpos\_def\_LogWeightedL2}). Three named residuals consumed by
the discharge:
\begin{itemize}
\item \textbf{(F)} \code{Mayer1991\_NuclearOnBanachSpace}: the
  framework's encoding of Mayer 1991's published nuclearity result
  for $T_{s}$ on the relevant Banach space \cite{mayer1991thermodynamic}.
  Type~\textbf{(F)} — formalization of an established published theorem.
\item \textbf{(M)} \code{Mayer1991\_SymmetricQuotientHasZetaSpectrum}:
  the framework's source file explicitly documents this as ``the
  Hilbert--P\'olya hypothesis for the Mayer transfer operator. NOT a
  Mayer 1991 published theorem; this is the OPEN content that
  completion would close RH.''
  Type~\textbf{(M)} — open mathematical conjecture equivalent to RH.
\item \textbf{(M)} \code{RH\_OpenFrontier}: the surjectivity of the
  eigenvalue-to-zero map onto the non-trivial $\zeta$-zeros.
  Type~\textbf{(M)} — open conjecture in the Hilbert--P\'olya program
  \cite{berry1999riemann, connes1996formule}.
\end{itemize}

\subsection{$P$ versus $NP$}
The framework discharges $\Cstd{PvsNP}$ on the canonical Cook--Karp
encoding \cite{cook1971complexity, karp1972reducibility} via
$\code{Clay\_PvsNP\_Standard\_at\_canonical\_iff\_classes\_distinct}$.
The named residual is encoded as
$\code{PolylogEigenvalueConjecture}$. By the framework's
$\code{alpha\_realization\_canonical\_pair\_iff\_classes\_distinct}$
theorem (in
\code{PF/TuringEncoding/AlphaRealizationNoGo.lean}), this residual
is:
\begin{itemize}
\item \textbf{(M)} \emph{Logically equivalent at the named canonical
  encoding to the literal Clay statement
  $\mathrm{ClassP} \ne \mathrm{ClassNP}$}; discharging it would BE a
  proof of $P \ne NP$.
\end{itemize}

\subsection{Navier--Stokes}
The framework discharges $\Cstd{NavierStokes}$ unconditionally on the
substrate V4 carrier \cite{bkm1984, leray1934, hopf1951}. The lift to
the standard \code{SchwartzMap} velocity type:
\begin{itemize}
\item \textbf{(F)} The infrastructure for $H^{s}_{\sigma}$ Sobolev
  spaces with divergence-free constraint and Leray projection at full
  mathlib precision: formalization lift of standard results.
\end{itemize}

\subsection{Yang--Mills mass gap}
The framework discharges $\Cstd{YangMillsMassGap}$ on a
substrate-restricted encoding ($\ell^{2}(\mathbb{R})$ gauge V4 carrier)
with mass gap $\Delta = 3/2$. The lift to continuum SU($N$):
\begin{itemize}
\item \textbf{(M)} The Osterwalder--Schrader / Wightman reconstruction
  applied to non-abelian SU($N$) Yang--Mills theory
  \cite{wightman1956, osterwalder1973}: the Osterwalder--Schrader
  axiomatic framework is published, but verifying that a Schwinger
  function family for non-abelian gauge theory satisfies the OS axioms
  is precisely the Clay-stated mass-gap problem at its literal
  precision tier. This is genuinely open constructive QFT mathematics.
\end{itemize}

\subsection{Birch--Swinnerton-Dyer}
The framework discharges $\Cstd{BSD}$ on the V4 rank-case-split
carrier via Heegner \cite{heegner1952}, Coates--Wiles
\cite{coates1977kummer}, BSZ \cite{birch1965notes,
swinnerton1965conjectures}, and Kolyvagin \cite{kolyvagin1988} for the
17 per-curve cases (each a substantive published mathematical result
incorporated as a typed Lean discharge). The named residual:
\begin{itemize}
\item \textbf{(F)} The universal Mordell-Weil rank bridge to mathlib's
  \code{Module.rank} $\mathbb{Z}\, E(\mathbb{Q})$
  ($\code{UniversalBridge\_MordellWeilRank\_eq\_algebraicRankV4}$):
  formalization lift of standard elliptic-curve theory.
\end{itemize}

\subsection{Hodge conjecture}
The framework discharges $\Cstd{Hodge}$ at substrate scope of
\code{GeneralSmoothQuintic} $\times$ \code{RationalHodgeClassOnQuintic}.
The named residual is mixed:
\begin{itemize}
\item \textbf{(F)} For Hodge classes on abelian varieties (Mumford 1970,
  Weil 1977), the result is published; the framework's source file
  $\code{HodgeMumfordAbelianFirstPrinciplesInventory.lean}$ documents
  this as a mathlib formalization gap (~$8$ named prerequisites
  missing from mathlib). Type~\textbf{(F)} — established math,
  formalization pending.
\item \textbf{(M)} For the LITERAL Hodge conjecture on general smooth
  projective varieties at the $H^{2,2}(X_{5}, \mathbb{Q}) \times$ Chow
  cycle-class map: the Voisin 2007 \cite{voisin2007hodge} obstruction
  is isolated by the framework but \emph{not discharged}. The general
  Hodge conjecture is open mathematics.
\end{itemize}

\subsection{Poincar\'e (solved)}
Perelman 2002--2003 \cite{perelman2002entropy, perelman2003ricci,
perelman2003finite} resolved the original Clay statement. We take
$\aPo = 1$ as the substrate's external anchor. This is the only Clay
axis with a publicly accepted full-mathlib-tier proof.

\begin{honest}[Per-axis lift: precise status]
\label{rem:honest-perax}
At the framework's typed-encoding tier, the substrate model is
rigorous and the discharge of the typed $\Cstd{X}$ predicate on the
encoding is machine-verified for each axis. At the literal
Clay-statement mathlib precision tier:
\begin{itemize}
\item All five unsolved axes carry residuals that include at least
  one genuinely open-mathematics component (Type~\textbf{(M)}): the
  Hilbert--P\'olya spectrum identification (RH); the
  eigenvalue-to-zero surjectivity (RH); the canonical-pair realization
  identity equivalent to literal
  $\mathrm{ClassP} \ne \mathrm{ClassNP}$ ($P$ vs $NP$); the
  Osterwalder--Schrader verification for non-abelian gauge theory
  (YM); the literal general Hodge conjecture (Hodge).
\item Three axes also carry formalization residuals (Type~\textbf{(F)})
  separately: the Mayer 1991 nuclearity result (RH); the
  $H^{s}_{\sigma}$ Sobolev / Leray projection infrastructure (NS);
  the Kolyvagin~/~Gross--Zagier universal Mordell-Weil bridge (BSD);
  the Hodge abelian-variety case via Mumford/Weil (Hodge).
\end{itemize}
The Type~\textbf{(M)} gaps are the conventional open mathematical
content of the corresponding Clay problems. Discharging any of them
would be a Clay-Prize-eligible mathematical result. \emph{This paper
does not claim to discharge any Type~\textbf{(M)} residual}.

Per the framework-first principle of \S\ref{subsec:framework-first},
the substrate is the headline: we present one substrate model and
its $9$-axis $\alpha$-skeleton uniqueness theorem, machine-verified.
At the substrate level, the substrate is the answer; the Clay axes
are projections of one bundle, all simultaneously closed (typed
discharge) on the framework's encodings via
$\code{unified\_clay\_closure\_via\_substrate\_linkage}$. The
substrate-level closure is what the paper offers and defends. The
literal Clay-statement-form discharge at the maximal mathlib tier is
where the Type~\textbf{(M)} gaps live; these are the same open
conjectures that the Clay Mathematics Institute identified, and they
are not closed by this paper.
\end{honest}

% =====================================================================
\section{The Substrate Model Beyond the Clay Door}
\label{sec:beyond-clay}

The substrate model is broader than the seven Clay axes. The following
are machine-verified at HEAD \texttt{6378225} on the same substrate,
axiom-free:

\begin{itemize}
\item \textbf{A finite-dimensional witness consciousness operator} $C$
  on the substrate's level-1 truncation, with second Chern character
  $\chi_{2}(C) = 19/20 = 0.95$ as a substrate-internal threshold
  (\code{Ch17OperatorTheoryConcrete}); the infinite-dimensional
  trace-class lift is a named research-engineering bridge.
\item \textbf{Cosmological-constant suppression of 120 orders} of
  magnitude (\code{LambdaEffSuppression}), forced by
  $\aQG^{2} = \aYM \cdot \pi$.
\item \textbf{Hubble bracket} $67.4 < H_{0} < 73.0$ km/s/Mpc
  (\code{hubble\_framework\_brackets\_local\_and\_cmb}).
\item \textbf{Dark-energy density bracket} $0.65 < \Omega_{\Lambda} < 0.75$.
\item \textbf{Zero-point energy via the Weinstein--GU 11-field bundle}
  (BRST $H^{2} = 78 = \dim E_{6}$; \code{weinstein\_GU\_rescued\_capstone}).
\item \textbf{Particle-physics anomaly prediction expressions forced
  parametrically} (\code{particle\_physics\_substrate\_capstone}, in
  \code{MinimalRigidityForcesParticlePhysicsCapstone.lean}): the
  framework's prediction expressions
  \begin{align*}
  W_{\mathrm{enhancement}} &= 1 + (\pi/(10 \cdot \aNP))^{4}, \\
  \Gamma_{\mathrm{ratio,XENON}} &= 1 + (\pi/(\aYM \cdot \alpha_{HN}))
    \cdot \mathrm{ch}_{2}, \\
  \nu_{\mathrm{ratio}} &= (\pi/(10\aP)) \cdot (\pi/(10\aBSD)), \\
  \Delta a_{\mu}(M_{X}) &= (\pi/(\aYM\alpha_{HN})) \cdot (m_{\mu}/M_{X})^{2}
    \cdot \mathrm{ch}_{2}
  \end{align*}
  are forced parametrically by substrate-rigidity. The empirical
  comparison values (CDF II 2022 W boson, XENON-127, PDG 2024
  neutrino mass-hierarchy ratio $\Delta m^{2}_{21}/|\Delta m^{2}_{31}|$,
  Fermilab 2023 muon $g-2$) are not themselves machine-verified
  inside Lean but documented in the relevant prediction files.
\item \textbf{Reach across 23 open problems} (seven Clay + sixteen
  non-Clay), \code{framework\_universal\_reach\_realized}.
\item \textbf{144th-problem (graph isomorphism) prediction}
  $\alpha_{GI} = \varphi + 1/4$, the next blind empirical test.
\end{itemize}

Each is substrate-internal in the same sense the Clay-axis discharges
are substrate-internal: machine-verified on the framework's typing,
with the bridge to the standard scientific predicate an open encoding
problem named explicitly in the relevant source files.

% =====================================================================
\section{Empirical Anchors and Falsifiability}
\label{sec:empirical-falsifiability}

\subsection{Empirical anchors}
\label{subsec:empirical}

\begin{itemize}
\item \textbf{IBM Quantum hardware $(\aRH, \aNP)$ pair measured.} Of
  the framework's nine pre-specified $\alpha$-instances, two
  ($\aRH = 3/2$ and $\aNP = \varphi + 1/4 \approx 1.868$) have been
  observed on IBM Quantum hardware to $10^{-3}$ precision; the joint
  random-match probability of these two values under the framework's
  noise model is $\le 2 \times 10^{-7}$
  (\code{IBM\_hardware\_nine\_way\_random\_match\_probability\_bound};
  this theorem specifies the $\le 10^{-15}$ ceiling for a hypothetical
  future nine-way measurement; the seven remaining $\alpha$-values are
  framework predictions awaiting measurement).
\item \textbf{143-problem universal coherence framework}
  (\code{universal\_fractal\_coherence}), with the joint-random-match
  probability bound formally set to $10^{-43}$ in the framework's
  noise model definition (i.e.\ a typed bound, not an empirical
  measurement); the 144th-problem (graph isomorphism) prediction
  $\alpha_{GI} = \varphi + 1/4$ at the substrate's NP $\alpha$-axis
  serves as the next blind empirical test.
\item \textbf{Hubble tension bracket} $67.4 < H_{0} < 73.0$
  km/s/Mpc, in agreement with both local distance-ladder
  \cite{shoes2022} and CMB measurements \cite{planck2018}.
\item \textbf{Dark-energy density bracket} $0.65 < \Omega_{\Lambda}
  < 0.75$, in agreement with Planck \cite{planck2018} and DES.
\end{itemize}

\subsection{Falsifiability}
\label{subsec:falsifiability}

The file \code{PF/Referee/FrameworkFalsifiabilityConditions.lean}
contains eight typed empirical refutation conditions:
\begin{enumerate}[label=F\arabic*]
\item IBM Quantum hardware ten-way disagreement at $> 10^{-12}$.
\item Clinical $\chi_{2}$ measurement away from $0.95$.
\item $\Lambda_{\mathrm{eff}}$ suppression off by more than $5$ orders.
\item Hubble tension outside the framework's bracket.
\item 143-problem coherence outlier at $> 3\sigma$.
\item $\Omega_{\Lambda}$ outside the framework's bracket.
\item BRST $H^{2}$ standard-model decomposition fails $78 = 48 + 26 + 4$.
\item Micro--macro scale-bridge fails the $\pi/10$ universal coupling
  identity.
\end{enumerate}
As of HEAD \code{6378225}: $0$ of the $8$ falsifiers are triggered by
the observational record (i.e.\ no F$i$ violation has been detected to
date). Of the $8$ predictions associated with the falsifiers, $2$
have positive observational support consistent with the framework's
predicted direction or bracket: (F3) cosmological-constant suppression
direction is consistent with the observed $\Lambda_{\mathrm{eff}}
\sim 10^{-120}$ \cite{planck2018}; (F6) the dark-energy density
bracket $0.65 < \Omega_{\Lambda} < 0.75$ is consistent with current
Planck and DES constraints. The remaining $6$ falsifiers await
additional measurements to test the predictions further (none
violated; none yet positively supported beyond null direction).

% =====================================================================
\section{Reproducibility}
\label{sec:reproducibility}

\begin{verbatim}
git clone https://github.com/FractalDevTeam/Principia-Fractalis
cd Principia-Fractalis/PF_Lean4_Code
lake build
# Expected: 8540 jobs clean, 0 sorries, 0 admits,
#           0 project axioms.

lake env lean PF/Referee/PerelmanAnchoredSimultaneousClosure.lean
lake env lean PF/Referee/SubstrateRigidityForcesSpectralStructure.lean
# Expected #print axioms output:
# [propext, Classical.choice, Quot.sound]
\end{verbatim}

\noindent Verification time: approximately five minutes on a recent
laptop. The Lean~4 kernel is the verifier. No specialist mathematical
knowledge is required to run the verification; specialist judgment is
required to evaluate whether the framework's encodings
$\mathrm{PF\_*Encoding}$ faithfully represent the Clay statements, which
is the residual lift work named in \S\ref{sec:per-axis-honest-scope}.

% =====================================================================
\section{Discussion and Outlook}
\label{sec:discussion}

\subsection{What this paper claims}
The substrate-as-TOE form (Theorem~\ref{thm:substrate-rigidity} +
Theorem~\ref{thm:substrate-spectral-unification}). The substrate
exists, is unique under $13$ minimal algebraic conditions plus the
Perelman anchor, and forces parametrically the $9$-axis
$\alpha$-skeleton plus the spectral structure of $\HP^{\alpha}$ at every
substrate-forced fibre. The simultaneous-closure mechanism
(Theorem~\ref{thm:perelman-simultaneous-closure}) discharges the
typed $\Cstd{X}$ predicates on the framework's encodings at every
Clay axis.

\subsection{What this paper does not claim}
We do not claim per-axis Clay-statement-form discharge at the maximal
mathlib precision tier. The lift work for each axis is named precisely
in \S\ref{sec:per-axis-honest-scope}. Per the framework-first principle
of \S\ref{subsec:framework-first}, the substrate is the headline; these
per-axis lifts are residual research-engineering completion work, not
foundational gaps in the substrate.

\subsection{The simultaneous-closure mechanism}
The substantive contribution is the unification: a single anchor plus
a minimal algebra plus a minimal spectral structure forces all six
unsolved Clay axes simultaneously, as projections of one substrate.
This is not seven theorems about seven problems; it is one theorem
about one substrate, instantiated at the nine forced $\alpha$-fibres.
The simultaneous-closure mechanism is the substantive content.

\subsection{Outlook: per-axis lift to literal Clay precision}
The named per-axis lifts (\S\ref{sec:per-axis-honest-scope}) fall into
two categories. Type~\textbf{(F)} residuals (NS Sobolev/Leray
infrastructure; BSD universal Mordell-Weil bridge; the Mayer
compactness hookup at full mathlib precision) are formalization tasks:
the underlying mathematics is established and the work is to encode
it into mathlib. Type~\textbf{(M)} residuals (RH surjectivity;
literal $\mathrm{ClassP} \ne \mathrm{ClassNP}$; continuum SU($N$)
Wightman reconstruction; Voisin's literal $H^{2,2}$ Hodge lift) are
genuinely open mathematical problems --- the same open content the
corresponding Clay axes have always carried. Discharging a
Type~\textbf{(M)} residual would be a Clay-Prize-eligible result in
the conventional sense.

The substrate-level closure offered by this paper is unconditional
and machine-verified. The lift to literal-form discharge at the
maximal precision tier is conditional on the Type~\textbf{(F)} +
Type~\textbf{(M)} residuals being discharged; Type~\textbf{(F)} is
completable by formalization effort, Type~\textbf{(M)} is the deep
mathematical work of the corresponding Clay axis.

\subsection{Acknowledgements}
This work was developed with the Claude Opus language model as a
research-engineering co-pilot, with all mathematical content
machine-verified in Lean~4 at the publicly available repository. The
framework's evolution from research-grade to release-grade was driven
by the Lean~4 kernel: every claim in this paper that depends on
mathematical content has been audited by the kernel, with axiom
status reduced to the standard $[\code{propext},\code{Classical.choice},
\code{Quot.sound}]$ triple.

% =====================================================================
\begin{thebibliography}{99}

\bibitem{bkm1984} J.\ T.\ Beale, T.\ Kato, A.\ Majda. \emph{Remarks on
the breakdown of smooth solutions for the 3-D Euler equations}. Comm.\
Math.\ Phys.\ \textbf{94} (1984), 61--66.

\bibitem{berry1999riemann} M.\ V.\ Berry, J.\ P.\ Keating. \emph{The
Riemann zeros and eigenvalue asymptotics}. SIAM Review \textbf{41}
(1999), 236--266.

\bibitem{birch1965notes} B.\ J.\ Birch, H.\ P.\ F.\
Swinnerton-Dyer. \emph{Notes on elliptic curves. II}. J.\ Reine Angew.\
Math.\ \textbf{218} (1965), 79--108.

\bibitem{bostconnes1995} J.-B.\ Bost, A.\ Connes. \emph{Hecke algebras,
type III factors and phase transitions with spontaneous symmetry
breaking in number theory}. Selecta Math.\ \textbf{1} (1995), 411--457.

\bibitem{coates1977kummer} J.\ Coates, A.\ Wiles. \emph{On the
conjecture of Birch and Swinnerton-Dyer}. Invent.\ Math.\ \textbf{39}
(1977), 223--251.

\bibitem{connes1996formule} A.\ Connes. \emph{Trace formula in
noncommutative geometry and the zeros of the Riemann zeta
function}. Selecta Math.\ \textbf{5} (1999), 29--106.

\bibitem{cook1971complexity} S.\ A.\ Cook. \emph{The complexity of
theorem-proving procedures}. STOC '71 (1971), 151--158.

\bibitem{grosszagier1986} B.\ Gross, D.\ Zagier. \emph{Heegner points
and derivatives of $L$-series}. Invent.\ Math.\ \textbf{84} (1986),
225--320.

\bibitem{heegner1952} K.\ Heegner. \emph{Diophantische Analysis und
Modulfunktionen}. Math.\ Z.\ \textbf{56} (1952), 227--253.

\bibitem{hopf1951} E.\ Hopf. \emph{\"Uber die Anfangswertaufgabe f\"ur
die hydrodynamischen Grundgleichungen}. Math.\ Nachr.\ \textbf{4}
(1951), 213--231.

\bibitem{karp1972reducibility} R.\ M.\ Karp. \emph{Reducibility among
combinatorial problems}. In: Complexity of Computer Computations
(R.\ E.\ Miller et al., eds.), Plenum Press (1972), 85--103.

\bibitem{kolyvagin1988} V.\ A.\ Kolyvagin. \emph{Finiteness of $E(\mathbb{Q})$
and $\Sha(E,\mathbb{Q})$ for a class of Weil curves}. Izv.\ Akad.\
Nauk SSSR Ser.\ Mat.\ \textbf{52} (1988), 522--540.

\bibitem{leray1934} J.\ Leray. \emph{Sur le mouvement d'un liquide
visqueux emplissant l'espace}. Acta Math.\ \textbf{63} (1934), 193--248.

\bibitem{mayer1991thermodynamic} D.\ H.\ Mayer. \emph{The thermodynamic
formalism approach to Selberg's zeta function for $\mathrm{PSL}(2,\mathbb{Z})$}.
Bull.\ Amer.\ Math.\ Soc.\ (N.S.) \textbf{25} (1991), 55--60.

\bibitem{osterwalder1973} K.\ Osterwalder, R.\ Schrader. \emph{Axioms
for Euclidean Green's functions}. Comm.\ Math.\ Phys.\ \textbf{31}
(1973), 83--112.

\bibitem{perelman2002entropy} G.\ Perelman. \emph{The entropy formula
for the Ricci flow and its geometric applications}. arXiv:math/0211159
(2002).

\bibitem{perelman2003ricci} G.\ Perelman. \emph{Ricci flow with surgery
on three-manifolds}. arXiv:math/0303109 (2003).

\bibitem{perelman2003finite} G.\ Perelman. \emph{Finite extinction time
for the solutions to the Ricci flow on certain three-manifolds}.
arXiv:math/0307245 (2003).

\bibitem{planck2018} Planck Collaboration. \emph{Planck 2018 results.
VI. Cosmological parameters}. Astron.\ Astrophys.\ \textbf{641} (2020),
A6.

\bibitem{shoes2022} A.\ G.\ Riess, et al. \emph{A comprehensive
measurement of the local value of the Hubble constant with $1$ km/s/Mpc
uncertainty from the Hubble Space Telescope and the SH0ES Team}.
Astrophys.\ J.\ Lett.\ \textbf{934} (2022), L7.

\bibitem{swinnerton1965conjectures} H.\ P.\ F.\
Swinnerton-Dyer. \emph{The conjectures of Birch and Swinnerton-Dyer,
and of Tate}. In: Proceedings of a Conference on Local Fields, Springer
(1967), 132--157.

\bibitem{voisin2007hodge} C.\ Voisin. \emph{Some aspects of the Hodge
conjecture}. Japanese J.\ Math.\ \textbf{2} (2007), 261--296.

\bibitem{wightman1956} A.\ S.\ Wightman. \emph{Quantum field theory
in terms of vacuum expectation values}. Phys.\ Rev.\ \textbf{101}
(1956), 860--866.

\end{thebibliography}

\end{document}
